\documentclass[11pt,a4paper]{article}
\usepackage[utf8]{inputenc}
\usepackage[T1]{fontenc}
\usepackage[french]{babel}
\usepackage{geometry}
\geometry{hmargin=2.5cm,vmargin=2.5cm}
\usepackage{enumitem}
\usepackage{amsmath}
\usepackage{amsfonts}
\usepackage{titlesec}
\usepackage{xcolor}

\titlelabel{\thetitle.\quad}
\titleformat{\section}{\large\bfseries\color{blue!70!black}}{}{0pt}{}

\title{\textbf{Synthèse des Automatismes en Mathématiques au Collège}}
\author{Progression de la 6\textsuperscript{e} à la 3\textsuperscript{e}}
\date{}

\begin{document}

\maketitle

\section{Classe de 6\textsuperscript{e} (Fin du Cycle 3)}
L'objectif est de consolider les acquis de l'école élémentaire, de ritualiser la fluidité du calcul et d'installer l'usage des nombres décimaux.

\subsection*{Nombres et calculs}
\begin{itemize}[label=\labelitemi,leftmargin=*]
    \item Maîtrise absolue des tables de multiplication et d'addition.
    \item Multiplier ou diviser un entier ou un décimal par $10$, $100$, $1000$ (mécanisme du glissement des rangs).
    \item Calcul mental sur des décimaux simples (ex: $2,5 \times 4$ ; $1,2 + 0,8$).
    \item Associer des fractions simples à leur écriture décimale ($\frac{1}{2} = 0,5$ ; $\frac{1}{4} = 0,25$ ; $\frac{3}{4} = 0,75$ ; $\frac{1}{10} = 0,1$).
    \item Compléter un nombre à $100$, à $1000$, ou à l'unité (ex: combien manque-t-il à $47$ pour atteindre $100$ ; combien manque-t-il à $0,3$ pour atteindre $1$).
    \item Trouver le double ou la moitié d'un nombre simple.
    \item Multiplier un nombre par $4$, $5$, $8$, $9$ ou $11$ en utilisant une astuce simple (ex: multiplier par $9$ revient à multiplier par $10$ puis à retrancher le nombre une fois).
    \item Ajouter ou soustraire un nombre qui se termine par $8$ ou $9$ en passant par la dizaine supérieure (ex: $47 + 9 = 47 + 10 - 1$).
    \item Calculer mentalement une expression simple en appliquant la priorité de la multiplication sur l'addition (ex: $3 + 4 \times 8$).
    \item Reconnaître si une écriture numérique est une somme ou un produit (ex: $2 \times 4 + 6$ est-il une somme ou un produit ?).
\end{itemize}

\subsection*{Grandeurs et mesures / Géométrie}
\begin{itemize}[label=\labelitemi,leftmargin=*]
    \item Convertir des longueurs et des masses (mètres, grammes et dérivés).
    \item Connaître les formules du périmètre d'un rectangle et du carré.
    \item Reconnaître et nommer les figures de base (triangle rectangle, rectangle, carré, cercle) à partir de leur description ou de leur codage.
    \item Calculer une durée simple (ex: une réunion commencée à 8h50 et terminée à 11h42).
    \item Reconnaître si une figure possède un axe de symétrie, en observant ou en pliant/superposant.
    \item Construire le symétrique d'un point ou d'une figure simple sur quadrillage.
    \item Reconnaître et nommer un angle droit, aigu ou obtus simplement en l'observant, sans le mesurer.
    \item Reconnaître, à partir de trois longueurs simples données, si un triangle peut être construit ou non (inégalité triangulaire).
\end{itemize}

\section{Classe de 5\textsuperscript{e} (Entrée dans le Cycle 4)}
Introduction des nombres relatifs, des premières notions de priorité opératoire complexe et début de l'initiation au calcul littéral.

\subsection*{Nombres et calculs}
\begin{itemize}[label=\labelitemi,leftmargin=*]
    \item Additionner et soustraire des nombres relatifs simples (ex: $-3 - 5 = -8$ ; $-2 + 7 = 5$).
    \item Appliquer les priorités opératoires de tête (ex: calculer $5 + 3 \times 2$ en prioritisant la multiplication).
    \item Simplifier des fractions simples par observation visuelle (ex: $\frac{10}{15} = \frac{2}{3}$ en repérant le diviseur commun 5).
    \item Calculer des pourcentages faciles de tête : $10\%$, $50\%$ ou $25\%$ d'une quantité (diviser par 10, par 2 ou par 4).
    \item Simplifier une écriture avec une lettre sans calcul (ex: $2 \times a$ s'écrit $2a$ ; $a \times b$ s'écrit $ab$).
    \item Reconnaître si une écriture littérale est une somme ou un produit (ex: $2(a+4)$ est-il une somme ou un produit ?).
    \item Calculer une vitesse, une distance ou une durée simple à partir de valeurs rondes.
    \item Calculer une quatrième proportionnelle simple (compléter un tableau de proportionnalité à trous).
\end{itemize}

\subsection*{Algèbre (Calcul littéral)}
\begin{itemize}[label=\labelitemi,leftmargin=*]
    \item Réduire une expression très simple (ex: $3x + 2x = 5x$ ; $a + a = 2a$).
\end{itemize}

\subsection*{Géométrie}
\begin{itemize}[label=\labelitemi,leftmargin=*]
    \item Identifier des droites parallèles et perpendiculaires d'après le codage d'une figure.
    \item Connaître la somme des angles d'un triangle ($180^\circ$) pour déterminer un angle manquant de tête.
    \item Reconnaître le symétrique d'un point ou d'une figure par symétrie centrale (par rapport à un point).
    \item Reconnaître un parallélogramme, un rectangle, un losange ou un carré grâce au codage de la figure.
    \item Reconnaître des angles alternes-internes ou correspondants égaux pour déterminer si deux droites sont parallèles.
    \item Reconnaître une bissectrice ou une médiatrice sur une figure codée.
    \item Calculer le périmètre ou l'aire d'un cercle (disque) à partir de sa formule, avec des valeurs simples.
    \item Reconnaître un triangle isocèle, équilatéral ou rectangle à partir d'un schéma codé.
    \item Calculer l'aire d'un carré, d'un rectangle ou d'un triangle à partir de sa formule, avec des valeurs simples.
\end{itemize}

\section{Classe de 4\textsuperscript{e} (Approfondissement du Cycle 4)}
Accélération de l'abstraction avec la multiplication des relatifs, les règles sur les puissances et le développement d'expressions littérales simples.

\subsection*{Nombres et calculs}
\begin{itemize}[label=\labelitemi,leftmargin=*]
    \item Multiplier et diviser des nombres relatifs (règle des signes, ex: $-4 \times (-5) = 20$).
    \item Manipuler les puissances de 10 (ex: $10^3 = 1000$ ; $10^{-2} = 0,01$).
    \item Passer d'un nombre décimal à sa notation scientifique simple, et inversement retrouver l'écriture décimale à partir de la notation scientifique.
    \item Calculer avec des fractions (effectuer mentalement $\frac{2}{3} \times \frac{5}{7}$ ou repérer un dénominateur commun évident comme $\frac{1}{2} + \frac{1}{4}$).
    \item Décomposer un petit nombre entier (inférieur à 30) en produit de facteurs premiers (ex: $60 = 2^2 \times 3 \times 5$).
    \item Reconnaître si la formule de distributivité simple peut s'appliquer à une écriture donnée, avant même de développer.
\end{itemize}

\subsection*{Algèbre (Calcul littéral)}
\begin{itemize}[label=\labelitemi,leftmargin=*]
    \item Développer une expression simple par distributivité simple : $3(x + 2) = 3x + 6$.
    \item Calculer la valeur d'une expression littérale pour une valeur de $x$ entière positive (ex: évaluer $2x + 5$ pour $x = 3$).
    \item Calculer la valeur d'une expression littérale contenant un carré pour une valeur de $x$ négative (ex: évaluer $x^2 + 3$ pour $x = -2$).
\end{itemize}

\subsection*{Géométrie}
\begin{itemize}[label=\labelitemi,leftmargin=*]
    \item Repérer l'hypoténuse d'un triangle rectangle immédiatement.
    \item Utiliser l'effet d'un agrandissement ou d'une réduction simple sur des longueurs (double, moitié).
    \item Écrire l'égalité de Pythagore dans un triangle rectangle, sans encore calculer.
    \item Prouver si un triangle est rectangle ou non à partir de ses trois longueurs.
    \item Calculer une longueur (l'hypoténuse ou un côté de l'angle droit) à l'aide de l'égalité de Pythagore, avec des valeurs simples.
    \item Reconnaître une translation.
    \item Reconnaître les hauteurs ou les médiatrices d'un triangle et savoir qu'elles sont concourantes.
    \item Reconnaître une configuration de Thalès (triangles emboîtés uniquement) et calculer une longueur manquante à l'aide de l'égalité des rapports, avec des valeurs simples.
    \item Utiliser la réciproque de Thalès pour prouver que deux droites sont parallèles ou non (triangles emboîtés uniquement).
\end{itemize}

\section{Classe de 3\textsuperscript{e} (Objectif DNB)}
Consolidation de tous les automatismes du collège pour être mobilisés instantanément lors de la première épreuve du Brevet.

\subsection*{Nombres et calculs}
\begin{itemize}[label=\labelitemi,leftmargin=*]
    \item Connaître par cœur les carrés parfaits des entiers de 1 à 12 ($1, 4, 9, 16, \dots, 144$).
    \item Appliquer instantanément les critères de divisibilité par 2, 3, 5, 9.
    \item Écrire un même nombre sous de multiples formes ($1,2 = \frac{12}{10} = \frac{6}{5} = 120\%$).
    \item Calculer $a\%$ d'un nombre pour $a = 1, 10, 25, 50, 100$.
\end{itemize}

\subsection*{Algèbre (Calcul littéral)}
\begin{itemize}[label=\labelitemi,leftmargin=*]
    \item Développer par distributivité simple et double.
    \item Factoriser une expression évidente (ex: $5x + 15 = 5(x + 3)$).
    \item Résoudre mentalement des équations du premier degré basiques (ex: $2x = 10 \implies x = 5$).
\end{itemize}

\subsection*{Organisation de données / Fonctions}
\begin{itemize}[label=\labelitemi,leftmargin=*]
    \item Calculer une moyenne simple ou déterminer une fréquence.
    \item Lire graphiquement l'image ou l'antécédent d'un nombre par une fonction.
\end{itemize}

\end{document}
